\section{Conclusions and Reflections}
\label{sec:conclusions}

\paragraph{What replicated.} The released code, after repairs, ran end to end on all eight splits and two stages. Followup F1 was 69.62 on FewRel (paper: 74.61) and 45.95 on WikiZSL (paper: 47.80). The method's broad performance level replicated. FewRel's improvement from feedback also replicated, although with lower recall than reported. WikiZSL's did not: feedback lowered F1 by 3.20. The released unlabelled pools (100,000 instances) do not match the paper's stated 10,000, a reproducibility detail that a reader would miss without counting the files.

\paragraph{How well it generalised.} My primary research question asked whether REPaL generalises to new relation types and domains. The evidence is limited and mixed. NYT reached 47.10 and SemEval between 19.64 and 30.99 with local generators. An archived GPT-4o run reached 35.67 on SemEval at the initial stage. The scientific-literature domain, which was the most novel test, remains untested because the annotation was not completed. So I can say that REPaL's performance varied substantially across the two existing datasets, but I cannot claim that it generalises to scientific text. My second question, about which definitions give the best synthesis, was only addressed indirectly through the SemEval analysis.

\paragraph{What I would do differently.} First, I would schedule annotation much earlier and recruit annotators at the start, because it was the bottleneck that blocked the arXiv evaluation, the agreement study and part of the judge design. Second, I would hold the generator fixed across all comparisons and reserve credits for it, because the mid-project move from GPT-4o to local models confounded the extension results. Third, I would run several seeds on every new dataset, because single-split results cannot separate a real effect from noise. Fourth, I would hold out unseen sentences from the unlabelled pool on new datasets, avoiding a fully transductive evaluation. Finally, I would check dataset properties such as pool size against the paper before the main runs, not after them, and would log library versions and GPU details in every run, since I could not recover them afterwards.
